%% Shared preamble for both submission.tex and editing.tex.
%% Edit packages here once; both wrappers \input it.

\usepackage{graphicx}
% True-to-size, left-justified figure whose caption width matches the image width.
% Measures the image into a save-box, then wraps image+caption in a minipage of that
% exact width. Use inside a figure/figure* environment:
%   \tcfig{figures/foo.pdf}{Caption text.}{fig:label}
\newsavebox{\tcfigbox}
\newcommand{\tcfig}[3]{%
  \sbox{\tcfigbox}{\includegraphics{#1}}%
  \noindent\begin{minipage}{\wd\tcfigbox}%
  \usebox{\tcfigbox}\par
  \caption{#2}\label{#3}%
  \end{minipage}%
}
% INTERIM fit-to-page variant: scales a figure to fit the print box (text width,
% capped at 160 mm tall, aspect preserved) -- use for figures not yet redrawn to true
% size. NOT WYSIWYG; switch to \tcfig once the source is resized to the print box.
\newcommand{\tcfigfit}[3]{%
  \noindent\includegraphics[width=\linewidth,height=160mm,keepaspectratio]{#1}\par
  \caption{#2}\label{#3}%
}
\usepackage{multirow}
\usepackage{amsmath,amssymb,amsfonts}
\usepackage{amsthm}
\usepackage{mathrsfs}

%% ---- House proof format for the Supplementary mathematical notes -------------------
%% One consistent, readable style to build every SI proof on. Rules:
%%   1. State each result in a `proposition` (or `lemma`) environment: BOLD head, ROMAN
%%      body (no italic statement), auto-numbered "Proposition 1, 2, ...".
%%   2. Prove in the amsthm `proof` environment: italic "Proof." head, roman body, QED box.
%%   3. Structure a multi-part proof with run-in bold \pfstep{...} headers -- NEVER an
%%      itemize/bullet list, and NEVER scattered \emph. Keep parentheticals to a minimum.
%%   4. Display every nontrivial equation (numbered). One idea per paragraph.
\newtheoremstyle{tcplain}% name
  {11pt plus 3pt minus 2pt}% space above -- clear gap so a Proposition separates from prose
  {11pt plus 3pt minus 2pt}% space below
  {}% body font: roman (empty = upright, not italic)
  {}% first-line indent
  {\bfseries}% head font: bold
  {.}% punctuation after head
  {\newline}% head followed by a line break -> the statement starts on its own line
  {}% head spec (default "Name Number")
\theoremstyle{tcplain}
\newtheorem{proposition}{Proposition}
\newtheorem{lemma}{Lemma}
% Run-in step header for structured proofs (replaces bullets and \emph run-ins).
\newcommand{\pfstep}[1]{\par\smallskip\noindent\textbf{#1.}\enspace}
% Run-in subsection header for Supplementary Notes: BOLD UPRIGHT (not bold-italic).
% Nature Portfolio uses bold upright for these run-in subheads; sn-jnl's \paragraph is
% \bfseries\itshape, so use \notesec instead of \paragraph inside the notes.
\newcommand{\notesec}[1]{\par\smallskip\noindent\textbf{#1.}\enspace}
%% -----------------------------------------------------------------------------------
%% SUPPLEMENTARY CROSS-REFERENCES -- Nature Portfolio style. VERIFIED 2026.07 against the
%% Nature Portfolio formatting guide (nature.com/documents/ncomms-formatting-instructions.pdf),
%% which lists the accepted and rejected forms explicitly:
%%     Correct:   "Supplementary Fig. 1"   "Supplementary Table 1"   "Supplementary Equation (1)"
%%     INCORRECT: "Fig. S1"   "Supplemental Table 1"   "See Supplementary Information"
%% So the word "Supplementary" does the disambiguating -- NOT an "S" glued to the number, and the
%% SI floats are numbered from 1 (see the numbering block at the top of the SI in backmatter.tex,
%% which also renames the float via \figurename/\tablename so CAPTIONS read
%% "Supplementary Figure 1 | ...", while \ref still yields the bare number).
%% Cite every SI item through these macros: they make the correct form the only convenient one.
\newcommand{\suppfig}[1]{Supplementary Fig.~\ref{#1}}
\newcommand{\supptab}[1]{Supplementary Table~\ref{#1}}
\newcommand{\suppnoteref}[1]{Supplementary Note~\ref{#1}}
%% Supplementary Notes are unnumbered \subsection*, so they need their own counter to become
%% \label/\ref-able. Hand-typed "Supplementary Note~5" had to be renumbered by hand every time a
%% Note was inserted; \suppnote removes that failure mode -- insert a Note anywhere and every
%% citation to every other Note renumbers itself.
\newcounter{suppnote}
\newcommand{\suppnote}[2]{%  #1 = label, #2 = title (may carry \secstatus)
  \clearpage
  \refstepcounter{suppnote}%
  \subsection*{Supplementary Note~\thesuppnote: #2}%
  \label{#1}%
}
%% -----------------------------------------------------------------------------------
\usepackage[title]{appendix}
\usepackage{xcolor}
\usepackage{pifont}% zapf dingbats for the editing-only status markers (\ding{51}/{55}/{110})
\usepackage{textcomp}
\usepackage{manyfoot}
\usepackage{booktabs}
\usepackage{placeins}% \FloatBarrier: keep figures inside their section (main figs before Methods)
\usepackage{algorithm}
\usepackage{algorithmicx}
\usepackage{algpseudocode}
\usepackage{listings}

\raggedbottom

% Drafting helper: figure placeholder box. Swap \figph{height}{desc} for
% \includegraphics[width=\linewidth]{figures/figN.pdf} once the draw.io export exists.
\newcommand{\figph}[2]{\fbox{\parbox[c][#1][c]{0.97\linewidth}{\centering #2}}}
